\chapter{آلة من عصبونات حية}
\chaptermark{آلة من عصبونات حية}
\label{sec:livingmachine}
\begin{chaptergoals}
\item لماذا تطلق مزرعة بلا قنوات بثّ رشقات شبكية، وما الذي يحدد الفاصل بينها؛
\item كيف تغيّر الشبكات الحية استجاباتها تحت التغذية الراجعة، وما الذي لم يثبت بعد؛
\item كيف تعمل العضيّة خزّانًا يُدرَّب قارئه الخطي في الخارج، في السيليكون؛
\item كيف عُلِّمت المزارع حتى الآن، ولماذا ما زال المعلّم خارج الطبق.
\end{chaptergoals}

\section{منصة اختبار بمخطط مكشوف}

في الفصل~\ref{sec:debugging} صحّحنا أخطاء آلة لا مخطط لها: يسمع المسبار ألف عصبون من أصل ثمانين مليارًا، وكل ما عدا ذلك علينا أن نخمّنه. والمهندس في هذا الموقف يفعل شيئًا بسيطًا: يجمّع قطعة صغيرة من الدارة على لوحة تجارب، حيث يرى كل مكوّن. ويمكن فعل ذلك مع العصبونات أيضًا.

تُفصل عصبونات القشرة بعضها عن بعض وتُزرع على شريحة فيها مصفوفة أقطاب، من بضع عشرات من الأقطاب إلى عشرات الآلاف. وفي غضون أسابيع قليلة تمدّ الخلايا زوائدها وتتصل في شبكة من آلاف العصبونات أو مئات الآلاف، أغلبها \nt{GLUT} ونسبة أصغر \nt{GABA} تتوقف على المستحضر: ففي مزارع الحُصين تبلغ نحو $6\%$ من العصبونات~\cite{Benson1994}. وكل قطب يستطيع أن يستمع، كمسبار الفصل~\ref{sec:debugging}، وأن يحقن تيارًا، كإشارة اختبار. والنوع الثاني من المنصات هو \emph{العضيّات} (\foreignlanguage{english}{organoids}): كُتل ثلاثية الأبعاد من النسيج العصبي قطرها نحو مليمتر، تُستنبت من خلايا جذعية بشرية~\cite{Lancaster2013}. على هذه المنصات تعمل الشبكات الحية شهورًا، وهناك مرافق يمكن فيها إجراء تجارب على العضيّات عن بُعد، عبر الشبكة، أكثر من مئة يوم متواصلة~\cite{Jordan2024}. ويسمى هذا المجال \emph{الحوسبة البيولوجية} أو \emph{ذكاء العضيّات}~\cite{Smirnova2023}.

تختلف المنصة عن الدماغ في أمرين مهمين. فهي ضئيلة: عصبوناتها أقل بمئة ألف مرة. وليس فيها قنوات بثّ: مزرعة القشرة لا تحوي عصبونات \nt{DOPA} ولا \nt{SERO} ولا \nt{NORA} ولا \nt{ACET}. إنها آلة لها ناقل بيانات وتحكم في التدفق، لكن لا سجلات تحكم لها. فلننظر ماذا تستطيع.

\section{ما تفعله الشبكة من تلقاء نفسها}

المزرعة المتروكة وشأنها لا تصمت، ولا تُصدر ضجيجًا منتظمًا. فمن حين لآخر تشتعل الشبكة كلها تقريبًا دفعة واحدة، في \emph{رشقة شبكية}، ثم تهدأ من جديد؛ والفواصل بين الرشقات الشبكية من ثانية إلى دقائق، بحسب المزرعة~\cite{Wagenaar2006}. والصورة مألوفة للمهندس: مضخِّم ذو تغذية راجعة موجبة قوية وبلا حِمل يتحول إلى مذبذب.

الآلية نعرفها. الروابط المتبادلة القوية بين عصبونات \nt{GLUT} تُطلق انهيارًا متسلسلًا، كما في الفصل~\ref{sec:gaba} عند الإخلال بشروط القضية~\ref{prop:ei}. والانهيار يستنزف بسرعة مخزون الناقل العصبي في المشابك~\eqref{eq:depression}، فتصمت الشبكة. ثم يتعافى المخزون بثابت زمني $\tau_{\text{rec}}$، وحين تتعافى النسبة $x^*$ التي تكفي لانهيار جديد تأتي رشقة شبكية جديدة. والفاصل بين الرشقات الشبكية
\begin{empheq}[box=\keybox]{equation}
T_{\text{burst}} \;\approx\; \tau_{\text{rec}}\,\ln\frac{1}{1-x^*} .
\label{eq:burstinterval}
\end{empheq}
مع $\tau_{\text{rec}}=0.8\,\text{s}$ من الفصل~\ref{sec:code} و$x^*=0.9$ نحصل على نحو ثانيتين، ومع $x^*=0.99$ على نحو أربع. هذا تقدير النموذج بقيمة $\tau_{\text{rec}}$ المعتمدة في الكتاب، وهو يقع عند الطرف القصير للفواصل المقيسة. أما القياس المباشر في المزارع الدقيقة فيُظهر أن مخزون الناقل يتعافى بعد الرشقة الشبكية في عشرات الثواني~\cite{CohenSegal2011}، أي أن $\tau_{\text{rec}}$ هناك أكبر؛ ومع قيمة كهذه تعطي الصيغة~\eqref{eq:burstinterval} عشرات الثواني والدقائق، كما في المزارع البطيئة. هذا مذبذب من نوع الاسترخاء، قريب لمولّد الإيقاع في الفصل~\ref{sec:muscles}: هناك كان يعيد شحنَه تعبُ مجموعات العصبونات، وهنا مخزونُ الناقل.

الرشقات الشبكية هي ما تفعله الشبكة حين لا يكون لديها ما تفعله. وفي الدماغ الحي تظهر صورة مشابهة في النوم العميق (الفصل~\ref{sec:sleep}) وفي الدماغ النامي، قبل أن تصله مداخل حقيقية. والتشابه يوحي بفكرة، لكن تطابق الآليات فرضية.

\section{كيف نعلّم الشبكة بلا معلّم دوباميني}

ما دام \nt{DOPA} غائبًا عن المزرعة، فعلينا نحن أن نقدّم إشارة ”أفضل أو أسوأ“، عبر الأقطاب. وتُظهر التجارب أن الشبكة على المنصة تغيّر استجاباتها فعلًا، وأطرف ما فيها هو \emph{بماذا} تُعلَّم.

\emph{معلّم يصمت.} تُحفَّز الشبكة عبر قطب واحد مرة كل ثانية إلى ثلاث ثوانٍ، إلى أن تظهر على قطب آخر الاستجابة المطلوبة بعد $50\pm10\ms$ من المنبه. وما إن تظهر حتى يُوقَف التحفيز خمس دقائق. وبعد عدة دورات كهذه تظهر الاستجابة المطلوبة على المنبه فورًا~\cite{Shahaf2001}. المكافأة هنا هي الصمت: المنبه يخلخل الشبكة، وإيقافه يتركها في الحالة التي أعطت الجواب الصحيح. هذا هو النزول على التدرج عبر اضطرابات عشوائية من الفصل~\ref{sec:dofa} (القضية~\ref{prop:gradient})، حيث يقوم الهزّ نفسه مقام الخطأ: نهزّ حتى يصير الجواب صحيحًا.

\emph{شبكة تلعب التنس.} في تجربة على منصة بمصفوفة أقطاب كثيفة، وُصل نحو مليون عصبون بلعبة تنس حاسوبية بمضرب واحد~\cite{Kagan2022}. كان موضع الكرة يُقدَّم تيارًا إلى أقطاب ”حسية“: أين الكرة، بترميز المكان؛ وكم تبعد، بالتردد. وكانت الفعالية في منطقتين ”حركيتين“ تحرّك المضرب إلى أعلى أو إلى أسفل. وكانت قاعدة التعلّم غير مألوفة. إذا صدّ المضرب الكرة تلقّت الشبكة منبهًا قصيرًا \emph{متوقَّعًا}؛ وإذا أخفق تلقّت بضع ثوانٍ من تيار عشوائي \emph{غير متوقَّع}. وفي عشرين دقيقة من اللعب طالت سلاسل الإصابات. لم يظهر تحسّن ذو دلالة داخل الجلسة إلا مع التغذية الراجعة الكاملة؛ ومع الصمت بدل المنبه لعبت المزرعة أيضًا في آخر الجلسة أفضل مما بلا تغذية راجعة، لكن بأثر أصغر، ولم يكن هناك تعلّم ثابت من جلسة إلى أخرى على مدى أيام. والتحليل المفصّل لهذه التجربة في الفصل~\ref{sec:dishbrain}.

فسّر الباحثون النتيجة بفرضية أن الشبكة تتصرف بحيث يصير دخلها متوقَّعًا~\cite{Friston2010}. وهي الفكرة نفسها التي رأيناها في الاقتصاد في النبضات في الفصل~\ref{sec:code} وفي التنبؤ بالإطار في الفصل~\ref{sec:recognition}: الآلة تنفق أقل حين يكون الدخل متوقَّعًا. لكن التجربة نفسها، وخاصة كلام الباحثين عن ”وعي“ المزرعة، لقيت نقدًا حادًا: الأثر صغير، ويمكن تفسيره بما هو أبسط~\cite{Balci2023}. والتقييم المنصف هو هذا: المزرعة على المنصة تغيّر سلوكها تحت تأثير التغذية الراجعة، وهذا مُقاس؛ أما أنها تتعلم تحديدًا التنبؤ بدخلها ففرضية.

\emph{شبكة تقود جسمًا.} وبطريقة مشابهة وُصلت مزرعة بـ”جسم“ افتراضي بسيط وعُلّمت أن تحرّكه نحو هدف، باختيار النمط المكاني الزماني لمنبهات التدريب~\cite{Bakkum2008}. لا دوبامين في أيٍّ من هذه التجارب: المعلّم نمط تيار يختاره المهندس. أما ما يتغير حين يصير المعلّم برنامجًا أو الدوبامين نفسه، ففي الأقسام~\babelsublr{\ref{sec:dishdopamine}--\ref{sec:cartpole}}.

\begin{figure}[t]
\centering
\begin{tikzpicture}[>=Stealth,
  blk/.style={draw,very thick,rounded corners=2pt,align=center,font=\footnotesize,minimum height=0.8cm,fill=blue!5}]
% (a) closed loop
\node[font=\small] at (1.5,4.3) {(أ) حلقة مغلقة};
\node[blk,text width=2.6cm] (G) at (1.5,3.3) {اللعبة: كرة ومضرب};
\draw[very thick,fill=orange!10,rounded corners=3pt] (0.5,0.2) rectangle (2.5,1.6);
\foreach \x in {0.7,1.0,...,2.35}{\foreach \y in {0.4,0.7,1.0,1.3}{\fill[gray!60] (\x,\y) circle (0.04);}}
\draw[->,thick,blue!60!black] (G.west) -| (-0.5,0.9) -- (0.5,0.9);
\node[font=\scriptsize,blue!60!black,anchor=east] at (-0.6,2.1) {الكرة $\leftarrow$ تيار};
\draw[->,thick,red!70!black] (2.5,0.9) -- (3.5,0.9) |- (G.east);
\node[font=\scriptsize,red!70!black,anchor=west] at (3.6,2.1) {النبضات $\leftarrow$ المضرب};
\node[font=\scriptsize] at (1.5,-0.2) {عصبونات على مصفوفة أقطاب};
\node[font=\scriptsize,align=center] at (1.5,-0.85) {إصابة: منبه متوقَّع\\إخفاق: تيار عشوائي};
% (b) reservoir
\begin{scope}[xshift=7.9cm]
\node[font=\small] at (1.6,4.3) {(ب) خزّان};
\node[blk,text width=1.1cm] (X) at (-0.4,1.0) {الدخل\\$u(t)$};
\draw[very thick,fill=orange!10] (1.6,1.0) circle (0.8);
\foreach \a/\r in {20/0.35,80/0.5,150/0.3,210/0.55,280/0.4,330/0.2,110/0.15}{\fill[red!60!black] ({1.6+\r*cos(\a)},{1.0+\r*sin(\a)}) circle (0.05);}
\node[font=\scriptsize] at (1.6,-0.2) {عضيّة: بلا معلّم};
\node[blk,text width=1.8cm,fill=green!8] (W) at (4.0,1.0) {القراءة\\$W_{\text{out}}$};
\draw[->,thick] (X) -- (0.8,1.0);
\draw[->,thick] (2.4,1.0) -- node[above,font=\scriptsize] {$\mathbf r(t)$} (W);
\draw[->,thick] (W) -- (4.0,2.3) node[above,font=\scriptsize] {$y(t)$};
\node[font=\scriptsize,green!40!black] at (4.0,0.3) {تتعلم بمعلّم};
\end{scope}
\end{tikzpicture}
\caption{طريقتان لجعل شبكة حية تحسب. (أ) حلقة مغلقة: موضع الكرة يُقدَّم تيارًا، ونبضات المنطقتين الحركيتين تحرّك المضرب، والمنبه المتوقَّع أو العشوائي بعد الضربة هو المعلّم. (ب) خزّان~\eqref{eq:reservoir}: لا تتلقى العضيّة إشارة تعليم، بل تفكّك الدخل إلى استجابات لاخطية كثيرة ذات ذاكرة؛ ولا يتعلم بالخطأ إلا قارئ خطي واحد في الخارج. وروابط العضيّة نفسها تتغير أيضًا في أثناء ذلك، بلا معلّم.}
\label{fig:livingmachine}
\end{figure}

\section{خزّان حي}

وهناك طريقة أخرى لاستعمال الشبكة الحية: ألّا نعلّمها أصلًا. تسمى هذه الطريقة في التقنية \emph{الحوسبة بالخزّان}~\cite{Maass2002,JaegerHaas2004}. نأخذ نظامًا لاخطيًا جاهزًا ذا ذاكرة، أيَّ نظام، بشرط أن تكون استجابته للدخل غنية وأن تتوقف على الماضي القريب. الدخل $u(t)$ يهزّه، فنحصل على متجه استجابات $\mathbf r(t)$، ولا يُدرَّب إلا القارئ الخطي
\begin{empheq}[box=\keybox]{equation}
y(t) \;=\; W_{\text{out}}\,\mathbf r(t) ,
\label{eq:reservoir}
\end{empheq}
الذي تُضبط أوزانه بقاعدة المربعات الصغرى~\eqref{eq:lms} من الفصل~\ref{sec:motorlearning}. يفعل الخزّان ما فعلته عصبونات \nt{GLUT} الصغيرة في الفصل~\ref{sec:motorlearning}: يفكّك الدخل إلى متجه طويل تصير فيه الأصناف المطلوبة قابلة للفصل بمستوٍ (الفصل~\ref{sec:dendrites}).

وقد صلحت عضيّة على مصفوفة أقطاب لدور هذا الخزّان. فأصوات الكلام، حين قُدّمت إليها أنماطًا من التيار، ميّزها قارئ خطي واحد بعد تدريبه بحسب المتكلم بدقة نحو $78\%$~\cite{Cai2023}. ودقة $78\%$ رقم يذكره الباحثون أنفسهم وتكرره الصحافة. النتيجة مفيدة، لكن من المهم أن نفهم أين ”الذكاء“ فيها: العضيّة تقدّم اللاخطية والذاكرة المتلاشية، أما التعلم بالخطأ فيجري في الخارج، في السيليكون (الشكل~\ref{fig:livingmachine}). ومع ذلك لا تبقى العضيّة على حالها: فبحسب الباحثين، تغيّرت في أثناء التدريب اتصاليتها الوظيفية نفسها، وهم يسمّون ذلك تعلّمًا بلا معلّم في العضيّة ذاتها~\cite{Cai2023}. ولم يُفصل بعدُ أيّ جزء من الدقة يعود إلى هذا التغير وأيّ جزء إلى القارئ.

\section{ما تقوله المنصة عن الآلة}

\emph{الطاقة.} بحسب التقدير~\eqref{eq:energy} تكلّف النبضة نحو $10^{-10}\,\text{J}$. فمزرعة من مليون عصبون بتردد نبضة واحدة في الثانية تنفق على نقل الإشارات
\begin{empheq}[box=\keybox]{equation}
P \;\approx\; 10^{6}\times1\,\text{s}^{-1}\times10^{-10}\,\text{J} \;=\; 10^{-4}\,\text{W},
\label{eq:dishpower}
\end{empheq}
أي عُشر ميلي واط. أما الإلكترونيات التي تحفّز هذه النبضات وتسجلها وتعالجها فتنفق أكثر بمراتب. وكما في الفصل~\ref{sec:debugging}، عنق الزجاجة ليس الحساب، بل الاتصال به.

\emph{القطع الناقصة.} تُظهر المنصة كم يستطيع ناقل \nt{GLUT} مع تحكم \nt{GABA} في التدفق، بلا سجلات تحكم: أن ينتظم ذاتيًا، ويولّد الرشقات الشبكية، ويغيّر استجاباته تحت تأثير التغذية الراجعة، ويكون خزّانًا جيدًا. وما لا يستطيعه بدونها هو أن يقرر بنفسه ما الذي يُعدّ نجاحًا. هذه الإشارة يقدّمها على المنصة المهندس، وفي الدماغ، كما رأينا في الفصول~\babelsublr{\ref{sec:dofa}--\ref{sec:acet}}، قنوات البثّ.

\emph{الأخلاقيات.} تُستنبت العضيّات من خلايا بشرية، وكلما ازدادت تعقيدًا صار السؤال أجدّ: هل يمكن لنسيج كهذا أن يشعر بشيء؟ والنظرة الهندسية توحي بجواب جزئي: الألم في الفصل~\ref{sec:pain} ليس إشارة من مستشعر واحد، بل منظومة إنذار كاملة ببوابات ومقبض لشدة الصوت وقنوات بثّ، وكل ذلك غائب عن المنصة. لكن هذا استدلال لا برهان، والمجال نفسه يقترح أن يجري التقييم الأخلاقي جنبًا إلى جنب مع التجارب، منذ البداية~\cite{Smirnova2023}.

\begin{keyblock}
\begin{proposition}[شبكة حية على المنصة]
\label{prop:livingmachine}
شبكة من عصبونات \nt{GLUT} و\nt{GABA} على مصفوفة أقطاب تولّد الرشقات الشبكية من تلقاء نفسها~\eqref{eq:burstinterval}، وتغيّر استجاباتها تحت تأثير التغذية الراجعة، ويمكن أن تكون خزّانًا~\eqref{eq:reservoir} لقارئ مدرَّب في الخارج. كل هذا مُقاس. أما أن هذه الشبكة تتعلم بقاعدة عامة ما، كأن تتنبأ بدخلها، ففرضية، والعبارات الرنانة عن ”وعيها“ لا يقبلها أغلب المختصين.
\end{proposition}
\end{keyblock}

\section{دوبامين بومضة ضوء}
\label{sec:dishdopamine}

التجربة المباشرة الوحيدة التي نعرفها، حيث أُعطيت المزرعة الدوبامين معلّمًا، أُجريت على منصة عضيّات يتصل بها الباحثون عن بُعد~\cite{Jordan2024}. أُضيف إلى السائل الذي يغمر العضيّات دوبامين داخل ”قفص“ حساس للضوء: فالجزيء المقيَّد لا يؤثر في المستقبلات. وتركيز الدوبامين المقيَّد $1\mM$. وحين يلزم أن تُكافأ الشبكة يُسلَّط عليها ضوء فوق بنفسجي: فينفتح القفص ويخرج الدوبامين إلى الحجم.

والحلقة مبنية هكذا. يُقدَّم المنبه نفسه مرة بعد مرة؛ فإذا أثار نبضات أكثر من ذي قبل أُطلقت الومضة وتحرر الدوبامين. وعلى مدى $240$ تكرارًا ازداد عدد النبضات المستثارة إجمالًا. ويكتب الباحثون أنفسهم أنه لا يمكن من تجربة واحدة كهذه أن يُستنتج أن الزيادة سببها الدوبامين~\cite{Jordan2024}.

هذه، في عين المهندس، أول محاولة لتجميع قاعدة العوامل الثلاثة~\eqref{eq:threefactor} في طبق: فعالية المرسل والمستقبِل تترك أثرًا، والإشارة العامة تقرر هل يُثبَّت التغيير. لكن الإشارة هنا موجبة فقط: المكافأة إما موجودة وإما غائبة، والجهاز لا يُصدر خطأً سالبًا $\delta<0$. ثم إن الدوبامين ينتشر ويُزال ببطء، كالتركيز في~\eqref{eq:lowpass}، فقد ترفع الومضات المتكررة مستواه العام لا غير. ولتمييز التعلم من ذلك تلزم تجربتان ضابطتان: العدد نفسه من الومضات في لحظات عشوائية لا صلة لها باستجابة الشبكة، والومضات نفسها بلا دوبامين مقيَّد. ويلزم أيضًا اختبار بعد الراحة: هل بقي التغيير بعد أن زال الدوبامين؟

\section{الدوبامين يعلّم السيليكون}
\label{sec:biohybrid}

في التجربة المعاكسة تعلّم الخلايا الحية الإلكترونيات~\cite{Keene2020}. تُستنبت خلايا تفرز الدوبامين في قناة دقيقة فوق عنصر عصبي الشكل عضوي، أي ترانزستور تؤدي موصلية قناته دور وزن المشبك. والدوبامين الخارج من الخلايا يتأكسد عند قطب العنصر، فيغيّر ذلك الموصلية. ويحاكي الجهاز أيضًا آلية إزالة الناقل من الشق، فيمكن أن يُغيَّر الوزن تغييرًا دائمًا وأن يُعاد أيضًا إلى ما كان عليه.

من حيث الدارة هذا مكامل تسريبي بدخل كيميائي: الوزن يراكم تدفق الدوبامين، و”الإزالة“ تحدد سرعة نسيانه، وهي دارة \foreignlanguage{english}{RC} نفسها التي في~\eqref{eq:lowpass}. والمهم هنا اتجاه الاتصال. فمسبار الفصل~\ref{sec:debugging} لا يرى سجلات البثّ لأنها كيميائية. وهذا العنصر يقرأ سجلًا كيميائيًا بالذات ويحوّله إلى وزن كهربائي. وهو، بالنسبة إلى واجهات الدماغ والحاسوب، طريق إلى مستشعر يسمع لا ناقل البيانات وحده، بل سجلات التحكم أيضًا.

\section{عضيّات بعصبوناتها الدوبامينية}
\label{sec:midbrainorg}

صار في الإمكان أن تُستنبت من الخلايا الجذعية البشرية عضيّات تشبه منطقة الدماغ التي توجد فيها عصبونات \nt{DOPA}. وفيها عصبونات دوبامينية عاملة تفرز الدوبامين~\cite{Jo2016}. وتُستنبت هذه العضيّات في الغالب لدراسة الأمراض. ولم نجد تجارب استُعمل فيها دوبامينها معلّمًا لشبكة أخرى.

وهذه، بالنسبة إلى الآلة المصنوعة من عصبونات حية، هي القطعة الناقصة. فمنصة الفصل~\ref{sec:livingmachine} ناقل بيانات وتحكم في التدفق بلا سجلات تحكم؛ وهنا مصدر إشارة البثّ موجود. ما ينقص هو الناقد: شبكة تحسب خطأ التنبؤ~\eqref{eq:ctd} وتتحكم في هذه العصبونات. ووصل عضيّة قشرية بعضيّة كهذه بحيث يُفرز الدوبامين بحسب الخطأ مسألة مفتوحة، وهو حتى الآن فرضية لا تجربة.

\section{عضيّة توازن عمودًا بلا دوبامين}
\label{sec:cartpole}

أكمل التجارب الحديثة يستغني عن الدوبامين تمامًا~\cite{Robbins2026}. وُصلت عضيّات قشرة الفأر في حلقة مغلقة بمسألة ”العمود على العربة“: فعالية العضيّة تحرّك العربة، وموضع العمود يعود إليها بالتحفيز. وبين المحاولات تتلقى العضيّة منبهات تعليم قصيرة عالية التردد. أيّها بالضبط، يختاره برنامج تعلّم معزَّز.

النتائج مقيسة:
\begin{itemize}
\item في أغلب العضيّات أعطت إشارات التعليم التي اختارها البرنامج نتيجة أفضل من إشارات مختارة عشوائيًا أو من غياب الإشارة؛
\item بعد $45$ دقيقة من الراحة لم يبقَ التحسن؛
\item حجب مستقبلات الغلوتامات بنوعيها، \foreignlanguage{english}{AMPA} و\foreignlanguage{english}{NMDA}، ألغى التحسن.
\end{itemize}

البند الأخير يدل على مكان ما تعلّمته العضيّة: في مشابك \nt{GLUT}، وعبر المستقبِل الذي يعمل في القسم~\ref{sec:weightstore} كاشفًا لتزامن الدخل والخرج. والبند الثاني يبيّن ما ينقص: التغير السريع موجود، والتثبيت غائب، كالمتعلم السريع في الفصل~\ref{sec:motorlearning} بلا المتعلم البطيء. أما دور المعلّم فتغيّر: حلّ برنامج محل المهندس الذي يختار نمط التيار، لكن المعلّم لا يزال خارج الطبق.

ولم نجد كذلك بين التجارب الأقدم على تعليم المزارع فوق مصفوفات الأقطاب تجربة واحدة استُعمل فيها الدوبامين: إشارة التعليم في كل مكان كانت تحفيزًا كهربائيًا.

\section{من يعلّم المزرعة}

\begin{table}[t]
\centering
\footnotesize
\setlength{\tabcolsep}{4pt}
\renewcommand{\arraystretch}{1.15}
\begin{tabular}{>{\raggedright}p{2.7cm}>{\raggedright}p{2.5cm}>{\raggedright}p{3.2cm}>{\raggedright\arraybackslash}p{4.4cm}}
\toprule
التجربة & مَن يعلّم & إشارة التعليم & ما هو معروف \\
\midrule
إيقاف التحفيز عند الاستجابة المطلوبة & المهندس & انتهاء التحفيز & تثبت الاستجابة، وهذا مُقاس \\
\foreignlanguage{english}{DishBrain} (الفصل~\ref{sec:dishbrain}) & المهندس & تيار متوقَّع أو عشوائي & نمو ذو دلالة في الجولات داخل الجلسة مع التغذية الراجعة الكاملة فقط، وأثر أصغر مع الصمت، وهذا مُقاس؛ غير ثابت من جلسة إلى أخرى؛ والسبب موضع نقاش \\
العمود على العربة & برنامج تعلّم معزَّز & منبهات عالية التردد يختارها البرنامج & أفضل من العشوائية، لكنه لا يبقى بعد الراحة، وهذا مُقاس \\
دوبامين بالومضة & قاعدة ”نبضات أكثر = مكافأة“ & دوبامين يحرّره الضوء & زيادة النبضات ملاحظة؛ دور الدوبامين غير مثبت \\
مشبك هجين حيوي & خلايا حية & دوبامين من الخلايا & تغيّر طويل لوزن العنصر وإعادته، وهذا مُقاس \\
عضيّات بعصبونات \nt{DOPA} & --- & --- & مصدر الدوبامين موجود؛ ولم يُستعمل معلّمًا \\
\bottomrule
\end{tabular}
\caption{من يعلّم العصبونات الحية على الشريحة، وبماذا.}
\label{tab:dishteachers}
\end{table}

في كل التجارب التي تتعلم فيها المزرعة نفسها، المعلّم حتى الآن في الخارج (الجدول~\ref{tab:dishteachers}): مهندس، أو برنامج، أو قاعدة وضعها مهندس. ولم يدخل الدوبامين الطبق إلا مرة واحدة، ودوره لم يُثبت بعد. أما تجميع قناة بثّ حقيقية في الطبق، بناقد وخطأ وأثر كما في الفصل~\ref{sec:dofa}، فهو الخطوة التالية التي لم يخطُها أحد بعد.


\section{خلاصة}

\begin{table}[t]
\centering
\footnotesize
\setlength{\tabcolsep}{4pt}
\renewcommand{\arraystretch}{1.15}
\begin{tabular}{>{\raggedright}p{2.8cm}>{\raggedright}p{4.2cm}>{\raggedright\arraybackslash}p{5.8cm}}
\toprule
التجربة & ما تفعله الشبكة الحية & ما هو معروف \\
\midrule
شبكة بلا دخل & رشقات شبكية بفواصل من ثانية إلى دقائق~\eqref{eq:burstinterval} & مُقاس؛ والآلية عبر استنزاف المشابك نموذج مقبول عمومًا \\
معلّم يصمت & تثبت الاستجابة المطلوبة حين يُوقَف التحفيز & مُقاس \\
لعبة التنس & تطول سلاسل الإصابات؛ والتحسن ذو الدلالة داخل الجلسة لا يظهر إلا مع التغذية الراجعة الكاملة & تغيّر السلوك مُقاس؛ التعلم من جلسة إلى أخرى غير ثابت؛ حجم الأثر وتفسيره بالتنبؤ بالدخل موضع خلاف \\
جسم افتراضي & حركة نحو الهدف بنمط منبهات مختار & مُقاس \\
خزّان & لاخطية وذاكرة من أجل القراءة~\eqref{eq:reservoir} & مُقاس؛ التعلم بالخطأ في الخارج، في السيليكون؛ وروابط العضيّة تتغير أيضًا \\
الطاقة & نحو $10^{-4}\,\text{W}$ لكل مليون عصبون~\eqref{eq:dishpower} & تقدير وفق~\eqref{eq:energy} \\
\bottomrule
\end{tabular}
\caption{ما أظهرته الآلات المصنوعة من عصبونات حية.}
\label{tab:livingmachine}
\end{table}

الآلة المصنوعة من عصبونات حية (الجدول~\ref{tab:livingmachine}) منصة اختبار نرى عليها ما يستطيعه ناقل البيانات والتحكم في التدفق بلا سجلات تحكم. يستطيعان الكثير، لكنهما لا يستطيعان أن يقررا بنفسيهما ما هو الجيد. على المنصة يؤدي هذا الدور المهندس بنمط التيار الذي يختاره؛ وفي الدماغ تؤديه أسلاك البثّ القليلة التي خُصّص لها وسط الكتاب.

\section{تمارين}

\begin{exercise}[\lvlA]
\label{ex:21:1}
\emph{الفاصل بين الرشقات الشبكية.} تستنزف الرشقة الشبكية المخزون حتى الصفر، ثم $\tau_{\text{rec}}\,\dd x/\dd t=1-x$، وتبدأ رشقة شبكية جديدة عند $x=x^*$؛ $\tau_{\text{rec}}=0.8\,\text{s}$. (أ)~أوجد الفاصل عند $x^*=0.9$ و$0.99$. (ب)~العتبة $x^*$ تتذبذب عشوائيًا بمقدار $\pm0.005$ حول $0.99$. في أي مدى تقع الفواصل؟
\end{exercise}

\begin{exercise}[\lvlA]
\label{ex:21:2}
\emph{طاقة المنصة.} (أ)~ما القدرة التي يعطيها التقدير~\eqref{eq:dishpower} للدماغ كله ($8.6\cdot10^{10}$ عصبون)؟ (ب)~تسجّل المنصة $1024$ قطبًا بتردد $20\,\text{kHz}$ و$10$ بتات للعينة. كم مرة يفوق نقلُ هذا التدفق بكلفة $1\,\text{nJ}$ للبت قدرةَ المزرعة نفسها؟
\end{exercise}

\begin{exercise}[\lvlB]
\label{ex:21:3}
\emph{تصميم: إخماد الرشقات الشبكية.} تحفيز متباعد عبر أقطاب كثيرة يجعل كل مشبك يطلق الناقل بتردد $r_b$؛ والمخزون $\tau_{\text{rec}}\,\dd x/\dd t=1-x-Ur_b\tau_{\text{rec}}x$، $U=0.5$، $\tau_{\text{rec}}=0.8\,\text{s}$. ما أصغر $r_b$ لا يبلغ عندها المخزون $x^*=0.9$؛ $0.99$، وكم يضعف المشبك عندئذ؟
\end{exercise}

\begin{exercise}[\lvlB]
\label{ex:21:4}
\emph{ذاكرة الخزّان لا تزيد على عدد العصبونات.} خزّان له $N$ مخرجًا $\mathbf r(t)$ يتلقى دخلًا أبيض $u(t)$ متوسطه صفر وتباينه واحد؛ و$m_k$ أكبر مربع لارتباط القراءة $\mathbf w_k\cdot\mathbf r(t)$ مع $u(t-k)$. برهن أن سعة الذاكرة $\mathrm{MC}=\sum_{k\ge0}m_k\le N$.
\end{exercise}

\begin{exercise}[\lvlB]
\label{ex:21:5}
\emph{محاكاة: خزّان في السيليكون.} الخزّان $\mathbf r(t+1)=\tanh(W\mathbf r(t)+\mathbf v\,u(t)+\mathbf c)$، $N=30$، و$W$ عشوائية نصف قطرها الطيفي $0.9$، $v_i\sim U(-0.5,0.5)$، $c_i\sim U(-0.2,0.2)$، $u\sim U(-1,1)$، والقراءة انحدار حَرْفي. أوجد سعة ذاكرة التمرين~\ref{ex:21:4} مع $\tanh$ وبدونها. أيّ الخزّانين يتذكر أكثر؟
\end{exercise}

\begin{exercise}[\lvlB]
\label{ex:21:6}
\emph{قراءة لخزّان حي.} تعطي العضيّة $1024$ قناة و$P=240$ تسجيلًا تدريبيًا من صنفين. (أ)~كم سمة $n$ يمكن إبقاؤها كي يكون التصنيف العشوائي قابلًا للفصل، بحسب صيغة التمرين~\ref{ex:19:3}، باحتمال لا يزيد على $1\%$؟ (ب)~بكم انحرافًا معياريًا تعلو دقة $78\%$ على $100$ تسجيل اختبار فوق التخمين؟
\end{exercise}

\begin{exercise}[\lvlC]
\label{ex:21:7}
\emph{بحث: معلّم بلا قناة بثّ.} اقترح بروتوكول تحفيز يطبّق على المنصة قاعدة العوامل الثلاثة من الفصل~\ref{sec:dofa} عبر $8\text{--}1024$ قطبًا، وتجربة ضابطة بقراءة مجمَّدة تميّز تعلّم الأوزان من انزياح حالة الشبكة. أيّ نتيجة تدحض فرضية التعلم؟
\end{exercise}
